\section{Arquitectura SaaS Multi-Tenant y App Móvil Kotlin}

\subsection{1. Objetivo General de la Plataforma}
Desarrollar y operar una plataforma SaaS (Software as a Service) multi-empresa que permita a diversos negocios y empresas gestionar, personalizar y publicar digitalmente su catálogo de productos y precios. Cada empresa cuenta con un código QR único que los clientes finales pueden escanear desde sus dispositivos móviles o mediante la aplicación nativa Android escrita en Kotlin.

\subsection{2. Jerarquía de Roles y Matriz de Control de Acceso}
La plataforma implementa un control de acceso basado en roles (RBAC) con aislamiento lógico estricto:

\begin{center}
\begin{tabular}{lp{5.5cm}p{6cm}}
\toprule
\textbf{Rol} & \textbf{Nivel de Jerarquía} & \textbf{Permisos y Alcance} \\
\midrule
\texttt{admin} & Administrador General SaaS & Acceso total a todas las empresas, gestión de usuarios globales, planes de suscripción y bitácora general de auditoría. \\
\texttt{empresa} & Panel Empresarial Multi-Tenant & Gestión exclusiva de los productos, categorías, precios e imagen corporativa de su propia empresa (\texttt{company\_id}). Aislamiento de datos estricto. \\
\texttt{cliente} & Usuario Final / Cliente Mobile & Consulta del catálogo público vía escaneo de código QR o consumo de APIs REST desde la aplicación móvil nativa Android en Kotlin. \\
\bottomrule
\end{tabular}
\end{center}

\begin{nota}[Principio de Aislamiento Multi-Tenant]
Ninguna empresa puede consultar, modificar o eliminar los productos o categorías pertenecientes a otra empresa. El backend valida en cada petición que el parámetro \texttt{company\_id} corresponda al tenant del usuario autenticado.
\end{nota}

\subsection{3. Modelo de Datos y Diagrama de Tablas MySQL}
El esquema relacional en MySQL está estructurado para soportar multi-tenancy eficiente y auditoría continua:

\begin{lstlisting}[language=SQL, caption={Estructura principal de tablas en MySQL}]
-- Tabla de Empresas (Tenants)
CREATE TABLE companies (
    id BIGINT AUTO_INCREMENT PRIMARY KEY,
    name VARCHAR(255) NOT NULL,
    slug VARCHAR(255) UNIQUE NOT NULL,
    email VARCHAR(255) UNIQUE NOT NULL,
    phone VARCHAR(50) NULL,
    plan ENUM('Basico', 'Plus', 'Premium') DEFAULT 'Basico',
    status ENUM('Activo', 'Inactivo') DEFAULT 'Activo',
    created_at TIMESTAMP, updated_at TIMESTAMP
);

-- Tabla de Usuarios con vinculo a Empresa
CREATE TABLE users (
    id BIGINT AUTO_INCREMENT PRIMARY KEY,
    name VARCHAR(255) NOT NULL,
    email VARCHAR(255) UNIQUE NOT NULL,
    password VARCHAR(255) NOT NULL,
    role ENUM('admin', 'empresa', 'cliente') DEFAULT 'cliente',
    company_id BIGINT NULL,
    FOREIGN KEY (company_id) REFERENCES companies(id) ON DELETE SET NULL
);

-- Tabla de Categorias por Empresa
CREATE TABLE categories (
    id BIGINT AUTO_INCREMENT PRIMARY KEY,
    company_id BIGINT NOT NULL,
    name VARCHAR(255) NOT NULL,
    description TEXT NULL,
    FOREIGN KEY (company_id) REFERENCES companies(id) ON DELETE CASCADE
);

-- Tabla de Productos y Precios
CREATE TABLE products (
    id BIGINT AUTO_INCREMENT PRIMARY KEY,
    company_id BIGINT NOT NULL,
    category_id BIGINT NULL,
    name VARCHAR(255) NOT NULL,
    description TEXT NULL,
    price DECIMAL(10,2) NOT NULL,
    is_active BOOLEAN DEFAULT TRUE,
    FOREIGN KEY (company_id) REFERENCES companies(id) ON DELETE CASCADE,
    FOREIGN KEY (category_id) REFERENCES categories(id) ON DELETE SET NULL
);
\end{lstlisting}

\subsection{4. Comparativa de Tecnologías Utilizadas}
\begin{longtable}{p{3.5cm}p{4.5cm}p{7.5cm}}
\toprule
\textbf{Tecnología} & \textbf{Componente} & \textbf{Propósito y Justificación Técnica} \\
\midrule
\endhead
\textbf{Laravel 11} & Backend Framework & Proporciona la arquitectura del servidor, enrutamiento web, controladores API REST y la capa de autenticación con Laravel Passport/Sanctum. \\
\textbf{MySQL 8.0} & Motor de Base de Datos & Almacena relacionalmente las empresas, productos, categorías, usuarios y registros inmutables de auditoría. \\
\textbf{Docker \& n8n} & Orquestación e IA & Servidor en contenedores que ejecuta el motor de automatización n8n para aplicar filtros inteligentes con IA sobre los productos. \\
\textbf{Kotlin (Android)} & Aplicación Móvil Nativa & Consume los endpoints JSON del catálogo y permite a los clientes escanear códigos QR con la cámara del celular. \\
\textbf{Laravel Blade} & Frontend Panel Web & Renderiza el panel de control administrativo responsivo con Tailwind CSS y contadores de sesión en tiempo real. \\
\bottomrule
\end{longtable}

\subsection{5. Flujo de Integración API con App Móvil Kotlin}
La comunicación entre la aplicación móvil Android y el servidor Laravel se realiza mediante peticiones HTTPS/HTTP con respuestas en formato JSON estructurado:

\begin{lstlisting}[language=bash, caption={Peticion HTTP enviada por la App Android Kotlin al escanear el QR}]
GET /api/v1/catalogo/tech-corp HTTP/1.1
Host: 127.0.0.1:8000
Accept: application/json
\end{lstlisting}

\begin{lstlisting}[language={}, caption={Respuesta JSON recibida por la App Kotlin}]
{
  "status": "success",
  "empresa": {
    "nombre": "Tech Corp Inc.",
    "slug": "tech-corp",
    "qr_url": "http://127.0.0.1:8000/q/tech-corp"
  },
  "categorias": [
    {
      "name": "Software & Licencias",
      "products": [
        { "id": 1, "name": "Licencia Anual LLM Chat Pro", "price": "4999.00" }
      ]
    }
  ]
}
\end{lstlisting}
